% 叩き台：解答の引数は空欄。問題本体は大問ごとに1ページ。
% レビュー反映：疑問詞はWhat単独に限定。疑問文には応答例を添え、長い和訳は別行に置く。
\MajorQuestionLesson{疑問詞What と一般動詞}
\TypeQuestion{会話文を書き写し、日本語に訳しなさい。}{generalverb_wh_copy_translate}
\SetRowSpace{4mm}
\EngCopyTranslate{What do you have?\enspace---\enspace I have a pen.}{}
\EngCopyTranslate{What do they study?\enspace---\enspace They study English.}{}
\EngCopyTranslate{What do you like?\enspace---\enspace I like music.}{}
\EngCopyTranslate{What do they use?\enspace---\enspace They use this room.}{}

\LessonPageBreak
\TypeQuestion{英文を書き写したあと、主語に下線を引き、動詞を丸で囲みなさい。\par{\small ※会話文では、両方の文に印を付けること。}}{generalverb_wh_copy_structure}
\SetRowSpace{2.5mm}
\EngWriteQuestion{\EngDialogue{What do you have?}{I have a pen.}}{}
\EngWriteQuestion{\EngDialogue{What do they read?}{They read books.}}{}
\EngWriteQuestion{\EngDialogue{Are you busy?}{No, I am not busy.}}{}
\EngWriteQuestion{\EngDialogue{What do you play?}{I play tennis.}}{}
\EngWriteQuestion{\EngSentence{They don't like music.}}{}
\EngWriteQuestion{\EngDialogue{What do you study?}{I study English.}}{}
\EngWriteQuestion{\EngDialogue{What do they use?}{They use this room.}}{}
\EngWriteQuestion{\EngSentence{They are busy.}}{}
\EngWriteQuestion{\EngDialogue{What is this?}{It is a pen.}}{}
\EngWriteQuestion{\EngDialogue{What do you read?}{I read books.}}{}

\LessonPageBreak
\TypeQuestion{日本語に合うように、誤りのある英文の必要なところだけ直して書きなさい。\par{\small ※短縮形の有無は原文に合わせること。右の返答は参考にすること。}}{generalverb_wh_error_correct}
\SetRowSpace{2mm}
\EngWriteQuestion{\EngDialogue{What you do have?}{I have a pen.}\par〔あなたは何を持っていますか。\enspace---\enspace 私はペンを持っています。〕}{}
\EngWriteQuestion{\EngDialogue{What are you study?}{I study English.}\par〔あなたは何を勉強しますか。\enspace---\enspace 私は英語を勉強します。〕}{}
\EngWriteQuestion{\EngDialogue{What they do like?}{They like music.}\par〔彼らは何が好きですか。\enspace---\enspace 彼らは音楽が好きです。〕}{}
\EngWriteQuestion{\EngDialogue{What you play?}{I play tennis.}\par〔あなたは何のスポーツをしますか。\enspace---\enspace 私はテニスをします。〕}{}
\EngWriteQuestion{\EngDialogue{What do read you?}{I read books.}\par〔あなたは何を読みますか。\enspace---\enspace 私は本を読みます。〕}{}
\EngWriteQuestion{\EngDialogue{What are this?}{It is a pen.}\par〔これは何ですか。\enspace---\enspace それはペンです。〕}{}
\EngWriteQuestion{\EngDialogue{What they use?}{They use this room.}\par〔彼らは何を使いますか。\enspace---\enspace 彼らはこの部屋を使います。〕}{}
\EngWriteQuestion{\EngDialogue{Are you play tennis?}{Yes, I do.}\par〔あなたはテニスをしますか。\enspace---\enspace はい、します。〕}{}

\LessonPageBreak
\TypeQuestion{日本語に合うように、適切な語句を丸で囲み、英文を書きなさい。\par{\small ※会話文では、疑問文と返答の両方を完成させて書くこと。}}{generalverb_wh_fill_choice}
\SetRowSpace{1.5mm}
\EngWriteQuestion{あなたは何を持っていますか。\enspace---\enspace 私はペンを持っています。\par\EngDialogue{What \EngChoice{do / are} you have?}{I \EngChoice{am / have} a pen.}}{}
\EngWriteQuestion{これは何ですか。\enspace---\enspace それはペンです。\par\EngDialogue{What \EngChoice{is / do} this?}{It \EngChoice{do / is} a pen.}}{}
\EngWriteQuestion{彼らは何を勉強しますか。\enspace---\enspace 彼らは英語を勉強します。\par\EngDialogue{What \EngChoice{are / do} they study?}{\EngChoice{I / They} study English.}}{}
\EngWriteQuestion{あなたは何のスポーツをしますか。\enspace---\enspace 私はテニスをします。\par\EngDialogue{What do you \EngChoice{play / am}?}{I \EngChoice{play / am} tennis.}}{}
\EngWriteQuestion{彼らは何を読みますか。\enspace---\enspace 彼らは本を読みます。\par\EngDialogue{What do they \EngChoice{are / read}?}{\EngChoice{They / I} read books.}}{}
\EngWriteQuestion{これらは何ですか。\enspace---\enspace それらは本です。\par\EngDialogue{What \EngChoice{is / are} these?}{They \EngChoice{are / do} books.}}{}
\EngWriteQuestion{あなたは何が好きですか。\enspace---\enspace 私は音楽が好きです。\par\EngDialogue{What do you \EngChoice{like / is}?}{I \EngChoice{am / like} music.}}{}
\EngWriteQuestion{彼らは何を使いますか。\enspace---\enspace 彼らはこの部屋を使います。\par\EngDialogue{What \EngChoice{do / is} they use?}{\EngChoice{I / They} use this room.}}{}

\LessonPageBreak
\TypeQuestion{疑問文に対する答えを一つ選び、その記号を丸で囲みなさい。\par 疑問文と選んだ返答を、続けて日本語に訳しなさい。}{generalverb_wh_response_choice}
\SetRowSpace{2mm}
\EngResponseChoice{What do you have?}{\mbox{（ア）\EngText{I am a student.}}\qquad\mbox{（イ）\EngText{I have a pen.}}\qquad\mbox{（ウ）\EngText{I play tennis.}}}{}
\EngResponseChoice{What do they study?}{\mbox{（ア）\EngText{They study English.}}\qquad\mbox{（イ）\EngText{I study English.}}\qquad\mbox{（ウ）\EngText{They are busy.}}}{}
\EngResponseChoice{What do you like?}{\mbox{（ア）\EngText{I am kind.}}\qquad\mbox{（イ）\EngText{They like music.}}\qquad\mbox{（ウ）\EngText{I like music.}}}{}
\EngResponseChoice{What do they read?}{\mbox{（ア）\EngText{I read books.}}\qquad\mbox{（イ）\EngText{They are students.}}\qquad\mbox{（ウ）\EngText{They read books.}}}{}
\EngResponseChoice{What is this?}{\mbox{（ア）\EngText{It is a pen.}}\qquad\mbox{（イ）\EngText{I have a pen.}}\qquad\mbox{（ウ）\EngText{They are pens.}}}{}
\EngResponseChoice{What do you play?}{\mbox{（ア）\EngText{I have a pen.}}\qquad\mbox{（イ）\EngText{I play tennis.}}\qquad\mbox{（ウ）\EngText{I read books.}}}{}
\EngResponseChoice{Do you like music?}{\mbox{（ア）\EngText{Yes, I am.}}\qquad\mbox{（イ）\EngText{Yes, I do.}}\qquad\mbox{（ウ）\EngText{Yes, you do.}}}{}
\EngResponseChoice{Are they students?}{\mbox{（ア）\EngText{No, they don't.}}\qquad\mbox{（イ）\EngText{No, I don't.}}\qquad\mbox{（ウ）\EngText{No, they aren't.}}}{}

\LessonPageBreak
\TypeQuestion{日本語に合う英文・会話文を、語句を並べ替えて作りなさい。\par{\small ※返答の語群はすべて使用。人名は日本語順、時の語句は文末に。}}{generalverb_wh_reorder}
\SetRowSpace{1.5mm}
\EngWriteQuestion{あなたは何を持っていますか。\enspace---\enspace 私はペンを持っています。\quad〔すべて使用〕\par\EngDialogue{\EngReorder{have / you / what / do}}{\EngReorder{a pen / have / I}}}{}
\EngWriteQuestion{彼らは何を読みますか。\enspace---\enspace 彼らは本を読みます。\quad〔すべて使用〕\par\EngDialogue{\EngReorder{read / what / they / do}}{\EngReorder{books / they / read}}}{}
\EngWriteQuestion{あなたは何を勉強しますか。\enspace---\enspace 私は英語を勉強します。\quad〔すべて使用〕\par\EngDialogue{\EngReorder{you / study / do / what}}{\EngReorder{English / I / study}}}{}
\EngWriteQuestion{あなたは何のスポーツをしますか。\enspace---\enspace 私はテニスをします。\quad〔すべて使用〕\par\EngDialogue{\EngReorder{play / do / what / you}}{\EngReorder{tennis / play / I}}}{}
\EngWriteQuestion{彼らは何を使いますか。\enspace---\enspace 彼らはこの部屋を使います。\quad〔すべて使用〕\par\EngDialogue{\EngReorder{do / they / use / what}}{\EngReorder{this room / use / they}}}{}
\EngWriteQuestion{これは何ですか。\enspace---\enspace それはペンです。\quad〔すべて使用〕\par\EngDialogue{\EngReorder{this / what / is}}{\EngReorder{a pen / it / is}}}{}
\EngWriteQuestion{彼らは何を持っていますか。\enspace---\enspace 彼らは本を持っています。\quad〔すべて使用〕\par\EngDialogue{\EngReorder{have / do / they / what}}{\EngReorder{books / have / they}}}{}
\EngWriteQuestion{あなたは何が好きですか。\enspace---\enspace 私は音楽が好きです。\quad〔すべて使用〕\par\EngDialogue{\EngReorder{like / you / what / do}}{\EngReorder{music / I / like}}}{}

\WordBank{疑問詞What と一般動詞}{
  \WB{I}{代名詞}{私は}{必ず大文字で書く}
  \WB{you}{代名詞}{あなた（たち）は}{}
  \WB{they}{代名詞}{彼ら・彼女ら・それらは}{}
  \WB{it}{代名詞}{それは}{}
  \WB{this}{代名詞等}{これ；この}{}
  \WB{these}{代名詞等}{これら；これらの}{\EngText{this}の複数形}
  \WB{a}{冠詞}{1つの／1人の}{訳さない場合もある}
  \WB{am / are / is}{be動詞}{～です／～にいます}{主語に合わせて使い分ける}
  \WB{do}{助動詞}{疑問・否定を作る}{単独では訳さない}
  \WB{not}{副詞}{～でない／～しない}{}
  \WB{don't}{短縮形}{しない；いいえ、しない}{\EngText{do not}の短縮形}
  \WB{aren't}{短縮形}{～でない}{\EngText{are not}の短縮形}
  \WB{yes}{応答語}{はい}{}
  \WB{no}{応答語}{いいえ}{}
  \WB{what}{疑問詞}{何；何の}{}
  \WB{play}{動詞}{（スポーツを）する}{}
  \WB{study}{動詞}{勉強する}{}
  \WB{like}{動詞}{～が好きである}{}
  \WB{read}{動詞}{読む}{}
  \WB{use}{動詞}{使う}{}
  \WB{have}{動詞}{持っている；食べる}{}
  \WB{English}{名詞}{英語}{最初の文字は大文字}
  \WB{music}{名詞}{音楽}{}
  \WB{tennis}{名詞}{テニス}{}
  \WB{book}{名詞}{本}{複数形 \EngText{books}}
  \WB{pen}{名詞}{ペン}{複数形 \EngText{pens}}
  \WB{room}{名詞}{部屋}{}
  \WB{student}{名詞}{生徒}{複数形 \EngText{students}}
  \WB{busy}{形容詞}{忙しい}{}
  \WB{kind}{形容詞}{親切な}{}
}

\PhraseBank{疑問詞What と一般動詞}{
  \PB*{What do you have?}{あなたは何を持っていますか。}
  \PB{I have a pen.}{私はペンを持っています。}
  \PB*{What do they study?}{彼らは何を勉強しますか。}
  \PB{They study English.}{彼らは英語を勉強します。}
  \PB*{What do you like?}{あなたは何が好きですか。}
  \PB{I like music.}{私は音楽が好きです。}
  \PB*{What do they use?}{彼らは何を使いますか。}
  \PB{They use this room.}{彼らはこの部屋を使います。}
}
